\section*{第 \ref{ch:g11:magnetic-fields} 章——磁现象与磁场}

\begin{solution}{exo:g11:magnetic-fields:1}
每一块都是一个完整的磁铁，有北极和南极：每个切面上
都会出现一个新磁极。新露出的两个面分别是 N 极和 S 极，因此
它们相互吸引——碎片会试图重新拼合。没有任何方法能孤立出一个磁极：
每次切割都会产生一对磁极。
\end{solution}

\begin{solution}{exo:g11:magnetic-fields:2}
(a) 错误：磁场在每一点只有一个方向。(b) 正确。
(c) 正确：每条磁感线都会穿过磁铁本身闭合。(d) 错误：
磁针会静止在\emph{沿}磁感线的方向上，即与之相切。
\end{solution}

\begin{solution}{exo:g11:magnetic-fields:3}
$B = \num{2e-7} \times 10 / 0.020 = \qty{1.0e-4}{T}$——是
地球磁场的两倍。
\end{solution}

\begin{solution}{exo:g11:magnetic-fields:4}
$n = 800/0.40 = \qty{2000}{m^{-1}}$;
$B = \num{1.26e-6} \times 2000 \times 1.5 \approx \qty{3.8}{mT}$.
$B \propto I$：电流加倍时，磁场变为 \qty{7.5}{mT}。
\end{solution}

\begin{solution}{exo:g11:magnetic-fields:5}
$F = ILB = 5.0 \times 0.10 \times 0.20 = \qty{0.10}{N}$——恰为质量
$m = F/g = 0.10/9.81 \approx \qty{10}{g}$ 的物体所受的重力。
\end{solution}

\begin{solution}{exo:g11:magnetic-fields:6}
\emph{1.} 她的航线指向真北偏西 $\ang{8}$：
$2000 \times \sin\ang{8} \approx \qty{2.8e2}{m}$，偏向西侧。

\emph{2.} 应沿磁针北向偏\emph{东} $\ang{8}$ 的方向前进。
\end{solution}

\begin{solution}{exo:g11:magnetic-fields:7}
$d = \mu_0 I / (2\pi B) = \num{2e-7} \times 20 / \num{5.0e-5}
= \qty{8.0}{cm}$。在距此类电缆一分米以内，指南针
指示的是导线的磁场而非地球磁场——应使其远离带电的导体。
\end{solution}

\begin{solution}{exo:g11:magnetic-fields:8}
$n = B/(\mu_0 I) = 0.010 / (\num{1.26e-6} \times 4.0)
\approx \qty{2.0e3}{m^{-1}}$；在 \qty{25}{cm}
$N = 1989 \times 0.25 \approx 500$ 匝。
\end{solution}

\begin{solution}{exo:g11:magnetic-fields:9}
$F = 10 \times 0.12 \times 0.50 = \qty{0.60}{N}$;
$a = F/m = 0.60/0.020 = \qty{30}{m/s^2}$——约为 $3g$：导体棒
确实会跳起来。
\end{solution}

\begin{solution}{exo:g11:magnetic-fields:10}
平衡时 $ILB = mg$，故
$I = \num{5.0e-3} \times 9.81 / (0.20 \times 0.10) \approx
\qty{2.5}{A}$。力必须竖直向上；$\vect B$ 水平且
与导线垂直，平展右手（掌心向上）可确定
所需的电流方向。
\end{solution}

\begin{solution}{exo:g11:magnetic-fields:11}
平展右手：(a) 指向页面顶部；(b) 指向右侧；
(c) 无力——导线与磁场平行。
\end{solution}

\begin{solution}{exo:g11:magnetic-fields:12}
\emph{1.} $B = \num{2e-7} \times 5.0 / 0.030 \approx
\qty{33}{\micro T}$；磁感线环绕导线，故正下方
处的磁场水平，与导线垂直，为东西方向。

\emph{2.} 磁针指向矢量和的方向：$\tan\theta = 33/20 =
1.67$，故 $\theta \approx \ang{59}$，偏离正北。
\end{solution}

\begin{solution}{exo:g11:magnetic-fields:13}
\emph{1.} $I = U/R = 12/2.0 = \qty{6.0}{A}$;
$n = 400/0.20 = \qty{2000}{m^{-1}}$;
$B_0 = \num{1.26e-6} \times 2000 \times 6.0 \approx \qty{15}{mT}$;
加入铁芯后，$100 B_0 \approx \qty{1.5}{T}$。

\emph{2.} $P = UI = 12 \times 6.0 = \qty{72}{W}$.

\emph{3.} 无电流则无磁场——而且软铁
不会保持磁性：负载会立即释放，这是设计所需。
\end{solution}

\begin{solution}{exo:g11:magnetic-fields:14}
\emph{1.} $F = NILB = 50 \times 2.0 \times 0.050 \times 0.30 =
\qty{1.5}{N}$，每边各受此力。两边的电流方向
相反，因此两力反向，构成使线圈转动的力偶。

\emph{2.} 导线与 $\vect B$ 平行时，不受拉普拉斯力。

\emph{3.} 转过半圈后，力偶将反向，
抵消原有的转动；换向器每半圈改变一次电流，
使线圈保持同一方向转动。
\end{solution}

\begin{solution}{exo:g11:magnetic-fields:15}
\emph{1.} $d = \num{2e-7} \times 100 / 1 = \qty{2.0e-5}{m} =
\qty{20}{\micro m}$——深入毫米厚铜导线的\emph{内部}。
单根导线附近的任何可到达点都达不到 \qty{1}{T}：强磁场
需要靠数千匝线圈叠加（即螺线管），再加上铁芯或
超导体。

\emph{2.} $\num{e8} / \num{5e-5} = \num{2e12}$：约相差十二个
数量级。
\end{solution}

\begin{solution}{pb:g11:magnetic-fields:1}
\textbf{1.} 指针本身是一块小磁铁；地球磁场
使其定向，北端沿 $\vect B$ 的水平分量方向，
即指向磁北。

\textbf{2.} $B = \sqrt{20^2 + 44^2} = \sqrt{2336} \approx
\qty{48}{\micro T}$，磁场斜向下，与水平面
夹角为 $\tan^{-1}(44/20) \approx \ang{66}$。

\textbf{3.} 航线为真北偏西 $\ang{6}$：
$30 \times \sin\ang{6} \approx \qty{3.1}{km}$，向西偏移。

\textbf{4.} 误差累加为 $\ang{10}$，则 $30 \times \sin\ang{10} \approx
\qty{5.2}{km}$——十度误差，偏离五千米。

\textbf{5.} 水平分量：靠近极点时，磁场
几乎竖直，水平分量趋于零，
指针便没有可对齐的方向了。

\textbf{6.} $n = 600/0.30 = \qty{2000}{m^{-1}}$.

\textbf{7.} $I = U/R = 12/1.5 = \qty{8.0}{A}$.

\textbf{8.} $B_0 = \mu_0 n I = \num{1.26e-6} \times 2000 \times 8.0
\approx \qty{20}{mT}$.

\textbf{9.} $50 \times \qty{20}{mT} = \qty{1.0}{T}$——是
钕铁硼磁铁磁极处磁场的两倍，且作用面积大得多。

\textbf{10.} $P = UI = 12 \times 8.0 = \qty{96}{W}$；在
\qty{10}{\minute} 内：$E = 96 \times 600 \approx \qty{58}{kJ}$。

\textbf{11.} 因为它能够释放：断开开关，磁场消失，
汽车恰好落在预定位置。永磁体则会一直吸着不放。

\textbf{12.} $F = NILB = 100 \times 3.0 \times 0.040 \times 0.25 =
\qty{3.0}{N}$.

\textbf{13.} 力偶：作用于轴两侧的
两个方向相反的力使线圈朝同一方向转动。大小相等的\emph{平行}
力只会使线圈侧移，无法使其转动。

\textbf{14.} 无：与磁场平行的导线不受
拉普拉斯力。

\textbf{15.} 每转半圈便使线圈中的电流反向，
正如力偶会改变符号一样——从而力矩始终驱使
线圈朝同一方向转动。

\textbf{16.} 将 $N$、$I$ 或 $B$ 任一加倍：这对力
$F = NILB$ 与每一个量均成正比。

\textbf{17.} $1.5 / \num{5.0e-5} = \num{3.0e4}$：相当于三万
个地球磁场。

\textbf{18.} $n = B/(\mu_0 I) = 1.5 / (\num{1.26e-6} \times 500)
\approx \qty{2.4e3}{m^{-1}}$——而其中每一匝都载有
\qty{500}{A} 的电流。

\textbf{19.} $P = RI^2 = 0.20 \times 500^2 = \qty{50}{kW}$——相当于
一个街区的取暖功率。实际绕组
为超导体：零电阻、零耗散，前提是
维持在高于绝对零度几度的温度。

\textbf{20.} 指南针 $\num{5e-5}$，起重机 $1$，MRI $1.5$：船只
位于两台机器之下五级，二者处于同一数量级 ——
介于星际空间的 $\num{e-10}$ 与中子星的 $\num{e8}$ 之间。唯一
改变的是电流（以及使之倍增的匝数、铁芯与超导体）。
同一磁场，三种用途。
\end{solution}
